\section{Homomorphisms}\label{S:sstd hom}

Let $\lambda$ be a partition of $n$ and $\mu$ be a composition of $n$. In \cite[Section~13]{GJ}, James constructed for each $\T \in \TT(\lambda,\mu)$ an $\F\sym{n}$-module homomorphism $\theta_\T \colon M_\F^{\lambda} \to M_\F^{\mu}$, and showed that $\{ \hat{\theta}_{\T} := \theta_{\T}|_{S_\F^{\lambda}}\colon \T \in \TTS(\lambda,\mu) \}$ is a basis for the space $\Hom_{\F\sym{n}}\big(S_\F^{\lambda},M_\F^{\mu}\big)$, unless $\Char(\F)=2$ and $\lambda$ is $2$-singular. In particular, this shows that every homomorphism in $\Hom_{\F\sym{n}}\big(S_\F^{\lambda},M_\F^{\mu}\big)$ is the restriction of a~homomorphism in $\Hom_{\F\sym{n}}\big(M_\F^{\lambda},M_\F^{\mu}\big)$ when $\Char(\F)\neq 2$.

In this section, we shall generalise these homomorphisms to obtain homomorphisms between~Specht modules and signed permutation modules. The next example shows that not every homomor\-phism in $\Hom_{\F\sym{n}}\big(S_\F^{\lambda},M_\F(\alpha|\beta)\big)$ is the restriction of a homomorphism in \linebreak $\Hom_{\F\sym{n}}\big(M_\F^{\lambda},M_\F(\alpha|\beta)\big)$, illustrating the difficulty of such generalisation.

\begin{Example}\label{Eg: James res}
Let $n\in \mathbb{Z}^+$ and let $\Char(\F)=p$ with $0 < p \leq n$. Let $\lambda=(1^n)$ and $(\alpha|\beta)=(\varnothing|(n))$. Then $M_\F^\lambda$ is the regular $\F\sym{n}$-module $\F\sym{n}$, while both $M_\F(\alpha|\beta)=M_\F(\varnothing|(n))$ and $S_\F^\lambda$ are isomorphic to the signature representation $\sgn$ of $\F\sym{n}$. Thus $\Hom_{\F\sym{n}}\big(S_\F^\lambda,M_\F(\alpha|\beta)\big)$ has dimension one.

On the other hand, $\Hom_{\F\sym{n}}\big(M_\F^\lambda,M_\F(\alpha|\beta)\big) \cong \Hom_{\F\sym{n}}(\F\sym{n},\sgn)$ has dimension $1$ with a~basis $\{\theta\colon \F\sym{n} \to \sgn \}$, where $\theta(1_{\sym{n}})=\epsilon$. For any $\sigma\in \sym{n}$, we have
\begin{gather*}
\theta(\sigma)=\sigma \cdot \theta(1_{\sym{n}})=\sgn(\sigma) \epsilon.
\end{gather*}
As a submodule of $M_\F^{\lambda} = \F\sym{n}$, the Specht module $S_\F^{\lambda}$ ($\cong \sgn$) is generated by $\sum\limits_{\sigma \in \sym{n}} \sgn(\sigma) \sigma$. We have
\begin{gather*}
\theta\left (\sum_{\sigma\in\sym{n}}\sgn(\sigma)\sigma\right )=\sum_{\sigma\in\sym{n}}\sgn(\sigma)\theta(\sigma)=\sum_{\sigma\in\sym{n}} \epsilon = (n!) \epsilon = 0,
\end{gather*}
since $p \leq n$. Thus $\theta|_{S_\F^\lambda} = 0$.

This example shows that the map $\Hom_{\F\sym{n}} \big(M_\F^{\lambda}, M_\F(\alpha|\beta)\big) \to \Hom_{\F\sym{n}} \big(S_\F^{\lambda}, M_\F(\alpha|\beta)\big)$ defined by $\theta \mapsto \theta|_{S_\F^{\lambda}}$ is not surjective in general.
\end{Example}

As such, to generalise the homomorphisms constructed by James for signed permutation modules, we should not attempt to generalise $\theta_\T$ and take its restriction to $S_\F^{\lambda}$, but have to generalise $\hat{\theta}_\T$ directly instead. In other words, we need to understand $\hat{\theta}_\T(e_\t)$.

\subsection{James's construction}
Let $\lambda$ be a partition of $n$ and $\mu$ be a composition of $n$. Let $\theta_{\T} \colon M_\F^{\lambda} \to M_\F^{\mu}$ be the $\F\sym{n}$-module homomorphism defined in \cite[Section~13]{GJ}. In this subsection, we study how $\theta_{\T}$ acts on the polytabloids in $M_\F^{\lambda}$.

Fix a $\lambda$-tableau $\t_0$, so that $\sym{n}$ acts on $\TT(\lambda,\mu)$ via $\t_0$, as described in Section~\ref{SS:Tableaux}.
Let $\T_0$ be the canonical $\lambda$-tableau of type $\mu$ associated to $\t_0$; recall that $\stab_{\sym{n}}(\T_0) = \sym{\mu}$.
Then each left coset of $\sym{\mu}$ in $\sym{n}$ corresponds to a $\lambda$-tableau of type $\mu$; let $d_\T\sym{\mu}$ be the left coset corresponding to $\T$.
Recall also the initial $\lambda$-tableau $\t^{\lambda}$, and let $d_{\t_0} = \t_0 \circ (\t^{\lambda})^{-1} \in \sym{n}$.
James defined $\theta_\T \in \Hom_{\F\sym{n}}\big(M_\F^{\lambda}, M_\F^{\mu}\big)$ so that
\begin{gather*}\theta_{\T}(d_{\t_0}\sym{\lambda}) = \sum_{\substack{d\sym{\mu} \in \sym{n}/\sym{\mu}: \\ d\sym{\mu} \subseteq R_{\t_0} d_\T \sym{\mu}}} d \sym{\mu}.\end{gather*}

Let $\t \in \TT(\lambda)$, and let $\rho_{\t} = \t \circ (\t_0)^{-1}$. Then $\rho_{\t} \cdot \t_0 = \rho_{\t} \circ \t_0 = \t$, so that $e_{\t} = e_{\rho_{\t} \cdot \t_0} = \rho_{\t} \cdot e_{\t_0}$. Thus
\begin{align*}\allowdisplaybreaks
\theta_{\T}(e_{\t}) &= \theta_{\T}(\rho_{\t} \cdot e_{\t_0})
= \theta_{\T}\left (\rho_{\t} \cdot \sum_{\sigma \in C_{\t_0}} \sgn(\sigma) (\sigma \cdot (d_{\t_0} \sym{\lambda}))\right ) \\
&= \sum_{\sigma \in C_{\t_0}} \sgn(\sigma) (\rho_{\t} \sigma \cdot \theta_{\T}(d_{\t_0} \sym{\lambda}))
= \sum_{\sigma \in C_{\t_0}} \sgn(\sigma) \left(\rho_{\t}\sigma \cdot \sum_{\substack{d\sym{\mu} \in \sym{n}/\sym{\mu} : \\ d\sym{\mu} \subseteq R_{\t_0} d_\T \sym{\mu}}} d \sym{\mu} \right) \\
&= \sum_{\sigma \in C_{\t_0}} \sgn(\sigma) \sum_{\substack{d\sym{\mu} \in \sym{n}/\sym{\mu} : \\ d\sym{\mu} \subseteq \rho_{\t} \sigma R_{\t_0} d_\T \sym{\mu}}} d \sym{\mu}
= \sum_{d \sym{\mu} \in \sym{n}/\sym{\mu}} a_{\rho_\t^{-1}d,d_\T} \, d\sym{\mu},
\end{align*}where
\begin{gather*}
a_{\rho_\t^{-1}d,d_\T}= \sum_{\substack{\sigma \in C_{\t_0} : \\ \sigma^{-1} \rho_{\t}^{-1} d \in R_{\t_0} d_{\T} \sym{\mu}}} \sgn(\sigma) =
\sum_{\substack{\sigma \in C_{\t_0} : \\ \sigma \rho_{\t}^{-1} d \in R_{\t_0} d_{\T} \sym{\mu}}} \sgn(\sigma).
\end{gather*}

We summarise this below.

\begin{Theorem} \label{T:JamesTheta}
Let $\lambda$ be a partition of $n$ and $\mu$ be a composition of $n$, and
let $\T \in \TT(\lambda,\mu)$.
The $\F\sym{n}$-module homomorphism $\hat{\theta}_{\T} \colon S_\F^{\lambda} \to M_\F^{\mu}$ constructed by James satisfies
\begin{gather*}
\hat{\theta}_{\T}(e_{\t}) = \sum_{d \sym{\mu} \in \sym{n}/\sym{\mu}} a_{\rho_\t^{-1}d,d_\T} \, d\sym{\mu},
\end{gather*} where
\begin{gather*}
a_{\rho_\t^{-1}d,d_\T} = \sum_{\substack{\sigma \in C_{\t_0} : \\ \sigma \rho_{\t}^{-1} d \in R_{\t_0} d_{\T} \sym{\mu}}} \sgn(\sigma).
\end{gather*}
\end{Theorem}

\subsection{Generalization of James's construction} \label{S:construction}
Fix a partition $\lambda$ of $n$ and a bicomposition $(\alpha|\beta)$ of $n$. In this subsection, we generalise James's construction of homomorphisms between Specht modules and Young permutation modules to obtain, for each $\dR$ in a subset $\R$ of $\sym{n}$ to be defined below (see Definition~\ref{D:RC}), a $\Z\sym{n}$-module homomorphism $\hJhom_\dR\colon S^\lambda_\Z \to M_\Z(\alpha|\beta)$ (see Theorem~\ref{T: hom}).


As before, we fix a $\lambda$-tableau $\t_0$, so that $\sym{n}$ acts on $\TT(\lambda, (\alpha|\beta))$ through $\t_0$, and denote the canonical $\lambda$-tableau of type $(\alpha|\beta)$ associated to $\t_0$ by $\T_0$. The following is the main theorem of this section:

\begin{Theorem}\label{T: hom}
	Let $\lambda$ be a partition of $n$, $(\alpha|\beta)$ be a bicomposition of $n$ and $\Gamma$ be a fixed left transversal of $\sym{\alpha|\beta}$ in $\sym{n}$. For each $\dR \in\R$, we have a $\Z\sym{n}$-module homomorphism $\hJhom_{\dR}\colon S_\Z^\lambda \to M_\Z(\alpha|\beta)$ given by
	\begin{gather*}
	\hJhom_{\dR}(e_{\t})=\sum_{d\in\Gamma} \a_{\rho_{\t}^{-1}d,\dR}\,(d \otimes \bb{1} \otimes \epsilon),\end{gather*}
	where \begin{gather*}\a_{\rho_{\t}^{-1}d, \dR} = \sum_{\substack{\sigma \in C_{\t_0} : \\ \sigma \rho_{\t}^{-1}d \in R_{\t_0} \dR \sym{\alpha|\beta} }} \sgn(\sigma)\varepsilon_\dR\big(\sigma \rho_{\t}^{-1}d\big),\end{gather*} and $\varepsilon_\dR\big(\tau \dR \xi_{\alpha}\xi_{\beta}^{+|\alpha|}\big) = \sgn(\xi_{\beta})$ for $\tau \in R_{\t_0}$ and $(\xi_{\alpha}, \xi_{\beta}) \in \sym{\alpha} \times \sym{\beta}$.
\end{Theorem}

By `reducing modulo $\Char(\F)$' the coefficients $\a_{\rho_{\t}^{-1}d, \dR}$ in $\hJhom_{\dR}(e_{\t})$, we obtain an $\F\sym{n}$-module homomorphism $\hJhom_{\dR}^\F\colon S^\lambda_\F\to M_\F(\alpha|\beta)$. Comparing this with $\hat{\theta}_{\T}$ in Theorem~\ref{T:JamesTheta}, one can see that the former is indeed a generalisation of the latter. The appearance of the map $\varepsilon_\dR$ in Theo\-rem~\ref{T: hom} also explains why our maps are indexed by elements of $\sym{n}$ instead of left cosets of $\sym{\alpha|\beta}$ (or equivalently, $\lambda$-tableaux of type $(\alpha|\beta)$), since $\varepsilon_\dR$ depends on $\dR$ and not on $\dR\sym{\alpha|\beta}$. A~little thought should convince the reader that to ensure that $\varepsilon_{\dR}$ is well-defined, $\dR$ may only run over a carefully chosen subset $\R$ of $\sym{n}$.

The remainder of this section is devoted to the proof of Theorem \ref{T: hom}.

\begin{Definition}\label{D:RC}
Define subsets of $\sym{n}$ as follows
\begin{gather*}
\R = \big\{ d \in \sym{n}\colon d^{-1} R_{\t_0}d \cap \sym{\alpha|\beta} \subseteq \sym{\alpha} \big\}, \\
\C = \big\{ d \in \sym{n}\colon d^{-1} C_{\t_0}d \cap \sym{\alpha|\beta} \subseteq \tlsym{\beta}{|\alpha|} \big\}.
\end{gather*}
\end{Definition}

These sets $\R$ and $\C$ of course depend on $\t_0$ and $(\alpha|\beta)$.

\begin{Lemma} \label{L:R}
Let $d \in \sym{n}$.
\begin{enumerate}\itemsep=0pt
\item[$(i)$] The following statements are equivalent:
\begin{enumerate}\itemsep=0pt
\item [$(a)$] $d \in \R$;
\item [$(b)$] $\stab_{R_{\t_0}} (\T_d) \subseteq d\sym{\alpha}d^{-1}$;
\item [$(c)$] whenever $\T_d(i,j) = \T_d(i,j')$ for some $(i,j),(i,j') \in [\lambda]$ with $j \ne j'$, we have $\T_d(i,j) = \bb{c}_k$ for some $k$.
\end{enumerate}
\item[$(ii)$] If $d \in \R$, then $\tau d \xi \in \R$ for all $\tau \in R_{\t_0}$ and $\xi \in \sym{\alpha|\beta}$.
\end{enumerate}
\end{Lemma}

\begin{proof} For part (i), firstly,
\begin{gather*}
\stab_{R_{\t_0}} (\T_d) = R_{\t_0} \cap \stab_{\sym{n}} (d \cdot \T_0) = R_{\t_0} \cap d \sym{\alpha|\beta} d^{-1}.
\end{gather*}
This proves the equivalence of (a) and (b).

Next, observe that a transposition $(a \;\; b)$ lies in $R_{\t_0}$ if and only if $a$ and $b$ label nodes in the same row of $\t_0$, i.e., $(\t_0)^{-1}(a) = (i,j)$ and $(\t_0)^{-1}(b) = (i,j')$ for some $i$. Furthermore, $(a \;\; b) \cdot \T_d = \T_d$ if and only if $\T_d((\t_0)^{-1}(a)) = \T_d((\t_0)^{-1}(b))$.
Since $\stab_{R_{\t_0}} (\T_d) = R_{\t_0} \cap d \sym{\alpha|\beta} d^{-1}$ is an intersection of conjugates of Young subgroups, it is generated by the transpositions it contains. Thus, $\stab_{R_{\t_0}} (\T_d)$ is generated by
\begin{gather*}
S = \big\{ (\t_0(i,j) \;\; \t_0(i,j')) \colon(i,j),(i,j') \in [\lambda],\, j \ne j',\, \T_d(i,j) = \T_d(i,j') \big\}.
\end{gather*}

Now,
\begin{gather*}
 (\t_0(i,j) \;\; \t_0(i,j')) \in d\sym{\alpha} d^{-1} \\
\quad {} \Leftrightarrow \big( d^{-1}(\t_0(i,j)) \;\; d^{-1}(\t_0(i,j'))\big) \in \sym{\alpha} \\
\quad {}\Leftrightarrow \exists\, (1 \leq k \leq \ell(\alpha)), \
\sum_{i=1}^{k-1} \alpha_i < d^{-1}(\t_0(i,j)), \ d^{-1}(\t_0(i,j')) \leq \sum_{i=1}^{k} \alpha_i\\
\quad {}\Leftrightarrow \exists\, (1 \leq k \leq \ell(\alpha)), \ \T_0 \big((\t_0)^{-1} \big(d^{-1}(\t_0(i,j))\big)\big) = \T_0 \big((\t_0)^{-1} \big( d^{-1}(\t_0(i,j'))\big)\big) = \bb{c}_k \\
\quad {} \Leftrightarrow \exists\, (1 \leq k \leq \ell(\alpha)), \ \T_d (i,j) = \T_d (i,j') = \bb{c}_k .
\end{gather*}
Hence (b) and (c) are equivalent.

For part (ii), observe that
\begin{gather*}
(\tau d \xi)^{-1} R_{\t_0} (\tau d \xi) \cap \sym{\alpha|\beta} = \xi^{-1}\big(d^{-1}R_{\t_0}d\cap \sym{\alpha|\beta}\big)\xi \subseteq
\xi^{-1}\sym{\alpha}\xi = \sym{\alpha}.\tag*{\qed}
\end{gather*}\renewcommand{\qed}{}
\end{proof}

We have analogous statements and proofs for $\C$ too.

\begin{Lemma} \label{L:C}
Let $d \in \sym{n}$.
\begin{enumerate}\itemsep=0pt
\item[$(i)$] The following statements are equivalent:
\begin{enumerate}\itemsep=0pt
\item [$(a)$] $d \in \C$;
\item [$(b)$] $\stab_{C_{\t_0}} (\T_d) \subseteq d\tlsym{\beta}{|\alpha|}d^{-1}$;
\item [$(c)$] whenever $\T_d(i,j) = \T_d(i',j)$ for some $(i,j),(i',j) \in [\lambda]$ with $i \ne i'$, we have $\T_d(i,j) = \bb{d}_k$ for some $k$.
\end{enumerate}
\item[$(ii)$] If $d \in \C$, then $\tau d \xi \in \C$ for all $\tau \in C_{\t_0}$ and $\xi \in \sym{\alpha|\beta}$.
\end{enumerate}
\end{Lemma}

Lemmas \ref{L:R} and \ref{L:C} give the following immediate corollary.

\begin{Corollary}\label{C:sstdRC}
If $d \in \sym{n}$ such that $\T_d \in \TTS(\lambda, (\alpha|\beta))$, then $d \in \R \cap \C$.
\end{Corollary}

In order to generalise the coefficient $a_{d,d_\T}$ in Theorem \ref{T:JamesTheta} to $\a_{d,\dR}$ in Definition \ref{D:Omega&a}, we need the following lemma which also explains the choice of the set $\R$.

\begin{Lemma} \label{L:projection}
Let $\dR \in \R$. There is a well-defined projection map $\pi_\dR\colon R_{\t_0}\dR\sym{\alpha|\beta} \to \sym{\beta}$ defined by $\tau \dR \xi_{\alpha} \tlxi{\beta}{|\alpha|} \mapsto \xi_{\beta}$ for all $\tau \in R_{\t_0}$ and $(\xi_\alpha,\xi_\beta) \in \sym{\alpha} \times \sym{\beta}$.
\end{Lemma}

\begin{proof}If $\tau \dR \xi_{\alpha}\tlxi{\beta}{|\alpha|} = \tau' \dR \xi'_{\alpha} \xi'^{+|\alpha|}_{\beta}$, then
\begin{gather*}\dR^{-1} R_{\t_0} \dR \ni \dR^{-1} \big(\tau^{-1} \tau'\big)\dR = \xi_{\alpha}\tlxi{\beta}{|\alpha|} \big(\xi'_{\alpha} \xi'^{+|\alpha|}_{\beta}\big)^{-1} \in \sym{\alpha|\beta},\end{gather*}
so that $\xi_{\alpha}\tlxi{\beta}{|\alpha|} \big(\xi'_{\alpha} \xi'^{+|\alpha|}_{\beta}\big)^{-1} \in \sym{\alpha}$ since $\dR \in \R$, forcing $\xi_{\beta} = \xi'_{\beta}$. The lemma thus follows.
\end{proof}

In view of Lemma \ref{L:projection}, we can make the following definition.

\begin{Definition}\label{D:Omega&a}
Let $d, \dR \in \sym{n}$ with $\dR \in \R$.
\begin{enumerate}\itemsep=0pt
\item [(i)] For $\omega \in R_{\t_0} \dR \sym{\alpha|\beta}$, let $\e_{\dR}(\omega) :=\sgn(\pi_{\dR}(\omega))\in \{\pm1\}$.
\item [(ii)] Let
\begin{gather*}
\Omega_{d,\dR} := \big\{ \sigma \in C_{\t_0}\colon \sigma d \in R_{\t_0} \dR \sym{\alpha|\beta} \big\}, \qquad
\a_{d,\dR} := \sum_{\sigma \in \Omega_{d,\dR} } \sgn(\sigma) \e_{\dR}(\sigma d) \in \Z.
\end{gather*} By convention, if $\Omega_{d,\dR}=\varnothing$ then $\a_{d,\dR}=0$.
\end{enumerate}
\end{Definition}

\begin{Remark} \label{Omega}
We give another description of $\Omega_{d,\dR}$ here. Restrict the left regular action of the symmetric group $\sym{n}$ on $\sym{n}/\sym{\alpha|\beta}$ to the subgroups $R_{\t_0}$ and $C_{\t_0}$, which partition $\sym{n}/\sym{\alpha|\beta}$ into $R_{\t_0}$-orbits and into $C_{\t_0}$-orbits respectively. Then $\sigma \in \Omega_{d,\dR}$ if and only if $\sigma \in C_{\t_0}$ and $\sigma \cdot (d\sym{\alpha|\beta}) \in R_{\t_0} \cdot (\dR\sym{\alpha|\beta})$. As such, $\Omega_{d,\dR} \ne \varnothing$ if and only if $C_{\t_0}d\sym{\alpha|\beta} \cap R_{\t_0} \dR\sym{\alpha|\beta} \ne \varnothing$. Furthermore, if $C_{\t_0}d\sym{\alpha|\beta} \cap R_{\t_0} \dR\sym{\alpha|\beta}$ contains precisely the distinct left cosets $d^{(1)}\sym{\alpha|\beta}, d^{(2)}\sym{\alpha|\beta}, \dotsc, d^{(r)}\sym{\alpha|\beta}$, then $\Omega_{d,\dR} = \bigcup\limits_{i=1}^r \Omega^{(i)}$ (disjoint union), where for each $i$, $\Omega^{(i)} = \{ \sigma \in C_{\t_0}\colon \sigma \cdot d\sym{\alpha|\beta} = d^{(i)}\sym{\alpha|\beta} \}$ and is therefore a left coset of $\stab_{C_{\t_0}} (d\sym{\alpha|\beta}) = C_{\t_0} \cap d\sym{\alpha|\beta} d^{-1}$. In particular, $\Omega_{d,\dR}$ is a union of some left cosets of $C_{\t_0} \cap d\sym{\alpha|\beta} d^{-1}$.
\end{Remark}

We collect together some important properties that $\a_{d,\dR}$ satisfies:

\begin{Lemma} \label{L:a}
Let $d,\dR \in \sym{n}$ with $\dR \in \R$.
\begin{enumerate}\itemsep=0pt
\item [$(i)$] If $\dR' = \tau \dR \xi_{\alpha}\tlxi{\beta}{|\alpha|}$ and $d' = \sigma d \eta_{\alpha}\eta_{\beta}^{+|\alpha|}$ for some $\tau\in R_{\t_0}$, $\sigma\in C_{\t_0}$ and $(\xi_\alpha,\xi_\beta),(\eta_\alpha,\eta_\beta) \in \sym{\alpha} \times \sym{\beta}$, then $\dR' \in \R$ and
 \begin{gather*}
 \a_{d',\dR'} = \sgn(\sigma)\sgn(\xi_\beta)\sgn(\eta_{\beta}) \a_{d,\dR}.
 \end{gather*}

\item [$(ii)$] If $C_{\t_0}d\sym{\alpha|\beta} \cap R_{\t_0} \dR \sym{\alpha|\beta} = \varnothing$ or $d \notin \C$, then $\a_{d,\dR} = 0$.
\item [$(iii)$] Suppose that $d \in \C$ and that $C_{\t_0}d\sym{\alpha|\beta} \cap R_{\t_0} \dR \sym{\alpha|\beta}$ is a disjoint union of $r$ left cosets of~$\sym{\alpha|\beta}$, with representatives $d^{(1)},\dotsc, d^{(r)}$. For each $i$, let $\sigma_i \in C_{\t_0}$ such that $\sigma_i d \in d^{(i)} \sym{\alpha|\beta}$, and let $\e^{(i)} = \sgn(\sigma_i) \e_{\dR}(\sigma_i d)$. Then
 \begin{gather*}
 \a_{d,\dR} = \big|C_{\t_0} \cap d \sym{\alpha|\beta} d^{-1}\big| \sum_{i=1}^r \e^{(i)}.
 \end{gather*}
\end{enumerate}
\end{Lemma}

\begin{proof} By Lemma \ref{L:R}(ii), $\dR' \in \R$. It is also easy to see that \begin{gather*}\Omega_{d',\dR'} = \Omega_{\sigma d,\dR}= \Omega_{d,\dR}\sigma^{-1}.\end{gather*} Take $\omega \in \Omega_{d',\dR'}$, say $\omega d'=\tau' \dR'\gamma_\alpha\gamma_\beta^{+|\alpha|}$ with $\tau'\in R_{\t_0}$ and $(\gamma_\alpha,\gamma_\beta) \in \sym{\alpha} \times \sym{\beta}$. Then \begin{gather*}\omega \sigma d=\tau' \tau \dR \xi_{\alpha}\gamma_\alpha\eta_{\alpha}^{-1}\big(\xi_\beta\gamma_\beta \eta_\beta^{-1}\big)^{+|\alpha|}\end{gather*} and hence we have $\e_{\dR'}(\omega d')=\sgn(\gamma_\beta)=\e_{\dR}(\omega \sigma d)\sgn(\eta_\beta)\sgn(\xi_\beta)$. Therefore{\samepage
\begin{align*}
\a_{d',\dR'} &= \sum_{\omega \in \Omega_{d',\dR'}} \sgn(\omega) \e_{\dR'}(\omega d')
= \sum_{\omega\sigma \in \Omega_{d,\dR}} \sgn(\sigma)\sgn(\omega\sigma) \e_{\dR}(\omega\sigma d)\sgn(\xi_{\beta})\sgn(\eta_{\beta}) \\
&= \sgn(\sigma)\sgn(\xi_{\beta})\sgn(\eta_{\beta}) \a_{d,\dR}.
\end{align*} This completes the proof of part (i).}

For part (ii), it is clear that $\Omega_{d,\dR} = \varnothing$ if $C_{\t_0}d\sym{\alpha|\beta} \cap R_{\t_0} \dR \sym{\alpha|\beta} = \varnothing$, so that $\a_{d,\dR} = 0$ in this instance. Next we assume that $d \notin \C$. By Lemma \ref{L:C}(i), there exist $(i,j), (i',j) \in [\lambda]$ with $i\ne i'$ such that $\T_d(i,j) = \bb{c}_k = \T_d(i',j)$ for some $k$. Let $\rho = (\t_0(i,j) \;\; \t_0(i',j))$. Then $\rho \in C_{\t_0} \cap d\sym{\alpha} d^{-1}$. If $\sigma \in \Omega_{d,\dR}$, say $\sigma d = \tau \dR \gamma_\alpha \gamma_\beta^{+|\alpha|}$ with $\tau \in R_{\t_0}$ and $(\gamma_\alpha,\gamma_\beta) \in \sym{\alpha} \times \sym{\beta}$, then
\begin{gather*}
R_{\t_0} \dR \sym{\alpha|\beta} \ni \tau \dR \gamma_{\alpha} \big(d^{-1}\rho d\big) \gamma_{\beta}^{+|\alpha|} = \sigma d \big(d^{-1}\rho d\big) = (\sigma \rho) d, \end{gather*}
so that $\sigma \rho \in \Omega_{d,\dR}$, and
\begin{gather*}
\sgn(\sigma \rho) \e_{\dR}(\sigma \rho d) = \sgn(\rho)\sgn(\sigma)\sgn(\gamma_\beta)= - \sgn(\sigma) \e_{\dR}(\sigma d).
\end{gather*}
As such, the contributions to the sum in $\a_{d,\dR}$ by $\sigma$ and $\sigma \rho$ cancel each other out. Consequently, $\a_{d,\dR} = 0$ and the proof of part (ii) is now complete.

For part (iii), for each $i$, let $\Omega^{(i)} = \{ \sigma \in \Omega_{d,\dR} \colon \sigma d \in d^{(i)}\sym{\alpha|\beta} \}$, so that $\Omega_{d,\dR}$ is a disjoint union of the $\Omega^{(i)}$'s (see Remark~\ref{Omega}). Fix $i$. There exist $\tau_i \in R_{\t_0}$ and $\big(\gamma^{(i)}_{\alpha},\gamma_{\beta}^{(i)}\big)\in\sym{\alpha}\times \sym{\beta}$ such that $\sigma_i d = \tau_i\dR \gamma_{\alpha}^{(i)} \tl{\big(\gamma_{\beta}^{(i)}\big)}{|\alpha|}$. For any $\omega \in \Omega^{(i)}$, we have
\begin{gather*} d^{-1} C_{\t_0} d \ni d^{-1} \big(\sigma_i^{-1}\omega\big) d = (\sigma_i d)^{-1} (\omega d) \in \sym{\alpha|\beta},\end{gather*}
so that $d^{-1} \big(\sigma_i^{-1}\omega\big) d \in \tlsym{\beta}{|\alpha|}$ since $d \in \C$. Thus
\begin{gather*}
\omega d = \sigma_i d \big(d^{-1} \sigma_i ^{-1} \omega d\big) = \tau_i \dR \gamma^{(i)}_{\alpha} \big(\gamma^{(i)}_{\beta}\big)^{+|\alpha|}\big( d^{-1} \sigma_i^{-1} \omega d\big),\end{gather*}
so that $\pi_{\dR}(\omega d) = \pi_{\dR}(\sigma_i d) d^{-1} \sigma_i^{-1} \omega d$ and hence
\begin{gather*}
\sgn(\omega) \e_{\dR}(\omega d) = \sgn(\omega) \sgn\big(\pi_{\dR}(\sigma_i d) d^{-1} \sigma_i^{-1} \omega d\big) = \sgn(\sigma_i)\e_{\dR}(\sigma_i d) = \e^{(i)}.\end{gather*}
Consequently,
\begin{align*}\allowdisplaybreaks
\a_{d,\dR} &= \sum_{\omega \in \Omega_{d,\dR}} \sgn(\omega) \e_{\dR}(\omega d)
= \sum_{i=1}^r \sum_{\omega \in \Omega^{(i)}} \sgn(\omega) \e_{\dR}(\omega d)
= \sum_{i=1}^r \sum_{\omega \in \Omega^{(i)}} \e^{(i)} \\
&= \sum_{i=1}^r |\Omega^{(i)}| \e^{(i)}
= \big|C_{\t_0} \cap d \sym{\alpha|\beta}d^{-1}\big| \sum_{i=1}^r \e^{(i)},
\end{align*} where the final equality is given by the following bijection: fix $\sigma^{(i)}\in \Omega^{(i)}$, we have a bijection between the sets $C_{\t_0} \cap d \sym{\alpha|\beta}d^{-1}$ and $\Omega^{(i)}$ given by $\sigma'\mapsto \sigma^{(i)}\sigma'$.
\end{proof}
